\section{Introduction}
\label{sec:introduction}

National aviation accounts must distinguish three concepts: who bought the fuel, which carrier
operated the aircraft and whether a flight stage was domestic or international. Petroleum
statistics commonly follow purchasers. IPCC inventories instead classify each flight from its
origin and destination, irrespective of airline nationality \citep{ipcc2006energy}. A national
airline can therefore buy fuel recorded under a domestic customer label while using much of it
on international services.

Pakistan makes the consequences visible. A movement-based calculation placed domestic
aviation near \SI{0.30}{\mega\tonne} \COe{} in FY2023--24, but had no fuel control and carried
a broad representation uncertainty. Public Oil Companies Advisory Council (OCAC) statistics
report JP-1 separately for PIA, other Pakistani-airline customers and foreign aviation. Read
literally, the resident-customer rows imply almost four times the movement estimate. That is
not evidence of fourfold domestic emissions: ``Domestic Avi--PIA'' identifies a purchaser, not
a domestic flight boundary.

The problem is a conserved partition rather than another uncertain percentage. Reassigning a
tonne from domestic aviation also changes international bunkers or another aviation account.
The source row must therefore remain intact while its destinations, method and unresolved
quantity stay visible. A second distinction follows. OECD estimates monthly territory- and
residence-based aviation emissions from flight data \citep{oecdaviation}. Those data can help
partition resident-customer fuel, but an estimate produced with that partition cannot then be
presented as an independent OECD validation. Gross fuel and flight accounts can still be
compared more independently, and their resident and foreign components can reveal whether an
apparently good total is internally coherent.

Aircraft-performance detail addresses a different layer. ICAO route factors and OpenAP
synthetic profiles can test the consequences of representative distance, route mix and
aircraft assignment \citep{icaoicec2024,sun2020openap}. Climate TRACE provides an additional
airport-level comparison, although surveillance and performance inputs may overlap with other
flight models \citep{climatetrace}. These sources strengthen the evidence hierarchy only when
their roles and shared inputs are declared in advance.

We therefore ask one question: can purchaser-based aviation-fuel statistics be reconciled with
flight-geography accounts without allowing aggregate agreement to conceal errors in the
underlying components? We answer it with nine OCAC fuel years, five complete fiscal-year OECD
comparisons, 42 matched months, Pakistan's official UNFCCC reporting and a hierarchy of
domestic activity models.

The paper makes three contributions. First, it introduces a source-row-conserving bridge that
separates purchaser, carrier residence and flight geography across domestic aviation,
international bunkers, military/government, general aviation, other uses and an explicit
residual. Second, it shows that near-closure of Pakistan's gross commercial aviation account
conceals persistent resident and foreign discrepancies of opposite sign. Third, it evaluates
domestic aviation through aggregate-movement, route-resolved and route--aircraft estimates,
with calibration, external comparison, structural scenarios and parametric uncertainty kept
distinct. Movement-to-cycle conversion, flight-geography classification and full-flight fuel
remain executable accounting safeguards rather than headline contributions.
